\subsection*{NT-C09--C13. Smaller class-group generators and exact computational input}

Let \(K\) be a number field of degree \(n\ge2\), discriminant
\(D=|\operatorname{disc}K|\), and \(r_1\) real places. All class
groups below are ordinary ideal class groups. Put
\[
 G=\operatorname{Cl}(K),\qquad
 G(T)=\langle[\mathfrak p]:N\mathfrak p\le T\rangle .
\]
The analytic input used for the uniform bounds is
\[
 {\cal R}_K:\quad
 \text{RH for }\zeta_{\mathbf Q}
 \text{ and for every }L_K(s,\chi),\quad \chi\in\widehat G,
                                                               \tag{C09.1}
\]
including \(\chi=1\), whose function is \(\zeta_K\).
Equivalently, it suffices to assume RH for \(\zeta_{\mathbf Q}\)
and \(\zeta_{H_K}\), where \(H_K\) is the ordinary Hilbert class
field: the all-prime abelian factorization NT-H11 has no
factor poles in the open critical strip. Thus \(M\) supplies
(C09.1). The designated input \(A\) suffices if \(H_K/\mathbf Q\)
is abelian; abelianness of \(K/\mathbf Q\) alone does not give
that conclusion.

\paragraph{NT-C09 and NT-C10: the precise published bounds.}
Greni\'e--Molteni, \emph{Breaking the 4 barrier for the bound of
a generating set of the class group}, Math.\ Comp.\ 94 (2025),
3145--3163, Theorems~1--4, author pp.~2--3, give the following
conclusions under (C09.1):
\[
 \begin{array}{ll}
 G(B_n)=G,&
 B_n=\left(4-\dfrac1{2n}+\dfrac1{2n^2}\right)(\log D)^2,\\[5pt]
 G(B'_n)\supseteq G^2,&
 B'_n=\left(4-\dfrac2{3n}+\dfrac3{8n^2}\right)(\log D)^2 .
 \end{array}                                                    \tag{C09.2}
\]
Without the respective \(n^{-2}\) corrections, these conclusions
hold for every \(n\ge5\). For \(n=2,3,4\), the verified ranges
are \(\log D\le l_-\) or \(\log D\ge l_+\), with
\[
 \begin{array}{c|rrr}
 n&2&3&4\\ \hline
 l_-&63&353&833\\
 l_+&12184&18455&27911\\ \hline
 l'_-&65&379&993\\
 l'_+&7968&9832&13108
 \end{array} .                                                  \tag{C09.3}
\]
The primed rows apply to \(G^2\).
The excluded middle intervals are intervals where the sharper
bound is not asserted, rather than a list of actual counterexamples.
For odd \(|G|\), squaring is an automorphism and both primed
bounds generate all of \(G\).

The proofs are in \S4, author pp.~10--12, and \S5,
pp.~12--13. Their analytic inputs are visible in \S2,
pp.~3--5 (class-character and Dedekind zero sums), and
Proposition~1, pp.~8--9 (RH for \(\zeta_{\mathbf Q}\)).
The small-discriminant input is their 2018
\emph{Explicit bounds for generators of the class group},
Theorem~3.5, equation~(3.10), author p.~6, using the
same class-character and Dedekind inputs.
The threshold computations, Odlyzko bounds and certified zero
tables remain credited parts of the published proof; they were
not numerically rerun here.

The containment in (C09.2) is the precise conclusion proved
at the start of the 2025 paper's \S5. Its introductory definition
of \(T'_K\) uses the phrase of generating \(G^2\); arbitrary
small prime classes need not themselves belong to \(G^2\).
We use the unambiguous subgroup-containment formulation throughout.

\paragraph{NT-C11: a finite certificate for a quotient, with exclusions.}
Here a smaller RH input suffices. Fix a subgroup \(H\le G\).
Assume RH only for the nontrivial class characters \(\chi\)
that are trivial on \(H\). No assumption on \(\zeta_K\) or
\(\zeta_{\mathbf Q}\) is needed in this paragraph.

Let \(L>0\), and let \(F:\mathbf R\to\mathbf R_{\ge0}\) be even, continuous,
compactly supported in \([-L,L]\), with nonnegative Fourier
transform. Assume the regularity needed in the Weil--Poitou
formula, for instance that \(F\) is piecewise \(C^1\) and
\((F(0)-F(x))/x\) is of bounded variation on \((0,\infty)\).
Define
\[
 \begin{split}
 I(F)&=\int_0^\infty
                  \frac{F(0)-F(x)}{2\sinh(x/2)}\,dx,\\
 J(F)&=\int_0^\infty\frac{F(x)}{2\cosh(x/2)}\,dx,\\
 W_F(\mathfrak p)&=\log N\mathfrak p
       \sum_{m\ge1}\frac{F(m\log N\mathfrak p)}
                              {(N\mathfrak p)^{m/2}},\\
 R_K(F)&=F(0)\{\log D-(\gamma+\log8\pi)n\}
                                      +nI(F)-r_1J(F).
 \end{split}                                                    \tag{C11.1}
\]
Let \(E\) be any set of prime ideals to be omitted. If
\[
       2\sum_{\mathfrak p}
             (-1)^{\mathbf1_E(\mathfrak p)}W_F(\mathfrak p)
                          >R_K(F),                            \tag{C11.2}
\]
then the classes of primes \(\mathfrak p\notin E\) with
\(N\mathfrak p<e^L\), together with \(H\), generate \(G\).
Only finitely many terms in (C11.2) are nonzero. This is a
quotient formulation of Greni\'e--Molteni (2025),
Theorem~5, author pp.~13--14, with the support stated in
the logarithmic variable in which the test function is evaluated.

\begin{proof}
Suppose the stated classes and \(H\) generate a proper
subgroup \(U<G\). A nontrivial character of \(G/U\), inflated
to \(G\), gives a character \(\chi\ne1\) trivial on \(H\)
and on every selected prime. Write
\[
       \phi(s)=\int_{\mathbf R}F(x)e^{(s-1/2)x}\,dx.
\]
The Weil--Poitou identity, in the normalization of the
2025 source (3), author p.~3, gives
\[
 2\sum_{\mathfrak p}\log N\mathfrak p
       \sum_{m\ge1}
       \frac{\operatorname{Re}(\chi(\mathfrak p)^m)
                          F(m\log N\mathfrak p)}
                              {(N\mathfrak p)^{m/2}}
       =R_K(F)-\sum_{\rho_\chi}\phi(\rho_\chi).
                                                               \tag{C11.3}
\]
The supplied RH hypothesis and Fourier positivity make every
zero term nonnegative. A prime outside \(E\) contributing to
the sum has norm below \(e^L\), hence character value one.
Indeed continuity and compact support force \(F(L)=0\),
so the endpoint creates no term. At a prime in \(E\),
\(\operatorname{Re}(\chi(\mathfrak p)^m)\ge-1\).
Since \(F\ge0\), the left side of (C11.3) is at least
the left side of (C11.2), whereas its right side is at most
\(R_K(F)\). This contradicts the strict inequality and
proves the assertion. The zero sum can equivalently be
read as the symmetric-height limit in the classical formula;
its terms are nonnegative here, so the comparison is valid.
\end{proof}

There is a directly computable sufficient version. Let \(T>1\),
\(L=\log T\), \(F(x)=(1-|x|/L)_+\), and write
\[
 {\mathfrak G}=\sum_{j\ge0}\frac{(-1)^j}{(2j+1)^2}.
\]
Then it suffices that
\[
 \begin{split}
 2\sum_{(N\mathfrak p)^m<T}
   (-1)^{\mathbf1_E(\mathfrak p)}
   \frac{\log N\mathfrak p}{(N\mathfrak p)^{m/2}}
        \left(1-\frac{m\log N\mathfrak p}{\log T}\right)
 >{}&\log D-n(\gamma+\log8\pi)-\frac{\pi r_1}{2}\\
    &+\frac{\pi^2 n/2+4{\mathfrak G}r_1}{\log T}.
 \end{split}                                                    \tag{C11.4}
\]
For \(H=1\) and \(E=\varnothing\), this is exactly the
Belabas--Diaz y Diaz--Friedman criterion, Corollary~2.2,
author pp.~6--7; its rounded Theorem~1.1 is on p.~3.
Here \(H=1\) denotes the trivial subgroup.

\begin{proof}
The Fourier transform of the triangular function is
\(L(\sin(tL/2)/(tL/2))^2\), with the value at zero taken
by continuity. It is nonnegative and meets all the
regularity assumptions. Also
\[
 1-F(x)\le x/L\quad(x\ge0),\qquad
 \int_0^\infty\frac{dx}{2\cosh(x/2)}=\frac\pi2.
\]
Thus
\[
 I(F)\le\frac{\pi^2}{2L},\qquad
 J(F)\ge\frac\pi2-\frac{4{\mathfrak G}}L.
 \]
For completeness, the constants follow by expanding
\[
 \frac1{2\sinh(x/2)}
       =\sum_{j\ge0}e^{-(j+1/2)x},\qquad
 \frac1{2\cosh(x/2)}
       =\sum_{j\ge0}(-1)^j e^{-(j+1/2)x}
\]
and integrating \(x\) times the series. The first integration
is monotone; the second is justified by the integrable
absolute sum of the \(x\)-weighted integrals. Their values
are \(\pi^2/2\) and \(4{\mathfrak G}\), respectively.
These inequalities bound \(R_K(F)\) above by the right
side of (C11.4); now apply (C11.2).
\end{proof}

\paragraph{NT-C12: what can actually be enumerated in polynomial time.}
Suppose an integral basis of the maximal order \(\mathcal O_K\)
and its full integer multiplication table are supplied.
Under (C09.1), one can deterministically enumerate, in
polynomial time in this input's bit length, prime ideals
whose classes generate \(G\), all of norm at most \(B_n+2\).
With the additional promise that \(|G|\) is odd, \(B'_n+2\)
may be used instead. The harmless additive two permits
certified rational rounding of the displayed logarithmic
bound; no unproved separation bound from an integer is needed.
The output is a list of integral ideal lattices.
This does not compute a relation lattice, the class number,
units, or a presentation of the class group.

\begin{proof}
First, the discriminant is computable by exact linear algebra
from the trace form. If all multiplication constants have
absolute value at most \(2^b\), the matrix of multiplication
by a basis element has that entry bound. Therefore each
trace-form entry has absolute value at most \(n^2 2^{2b}\).
Hadamard's inequality gives
\[
       \log_2 D\le 2bn+\frac{5n}{2}\log_2 n.                  \tag{C12.1}
\]
Consequently any integer \(T=O(1+(\log D)^2)\) is polynomial
in the supplied input length.

An integer \(T\) with \(B_n\le T\le B_n+2\) is obtainable
in polynomial time. Compute a rational upper approximation
to \(B_n\) with error less than \(1/2\), and take its ceiling.
For example, write \(D=2^m u\), \(1\le u<2\), and evaluate
\[
 \log u=2\sum_{j\ge0}
       \frac{1}{2j+1}\left(\frac{u-1}{u+1}\right)^{2j+1},
\]
as well as the same series for \(\log2\).
The ratios have absolute value at most \(1/3\); geometric
tail bounds give certified accuracy \(O((1+\log D)^{-1})\)
for \(\log D\), dividing the error allowance for \(\log2\)
by \(m+1\) before forming \(m\log2+\log u\).
This uses polynomially many rational bit operations. Squaring
and multiplying by the fixed rational coefficient gives
the required approximation. The same procedure applies
to \(B'_n\).

We describe an unconditional procedure to list all prime
ideals with norm at most this integer \(T\).
List rational primes \(p\le T\), even by trial division.
For each such \(p\), form the \(n\)-dimensional algebra
\[
                   A=\mathcal O_K/p\mathcal O_K
\]
over \(\mathbf F_p\), using the given table. Its Frobenius
\(F:a\mapsto a^p\) is \(\mathbf F_p\)-linear and is
computed by repeated squaring. Choose \(e\) with \(p^e\ge n\).
Then
\[
                   \operatorname{Nil}(A)=\ker F^e.           \tag{C12.2}
\]
Indeed a nilpotent in an \(n\)-dimensional algebra has
nilpotence index at most \(n\), whereas Frobenius is
injective on every finite residue field. Form the reduced
quotient \(R=A/\operatorname{Nil}(A)\). It is a product of
finite fields, and
\[
          B=\{b\in R:b^p=b\}\simeq\mathbf F_p^{\,g},
                                                               \tag{C12.3}
\]
where \(g\le n\) is the number of factors.

A basis of \(B\) is obtained by solving linear equations.
Use its elements successively to split \(R\) into the
simultaneous eigenspaces of their multiplication maps.
Every eigenvalue belongs to \(\mathbf F_p\), so it is enough
to scan \(a=0,\ldots,p-1\) and compute kernels of
\(m_b-aI\) on each current component. These maps are
diagonalizable because \(b^p=b\) and \(X^p-X\) has distinct
roots. The common eigenspaces are exactly the field factors:
a basis of \(\mathbf F_p^g\) separates every pair of
coordinates. All kernels are ideals, and the nonzero
resulting components give the maximal ideals of \(R\)
as kernels of the corresponding projections. Pulling
these back to \(A\) lists all maximal ideals \(W\subset A\).

If \(f=\dim_{\mathbf F_p}(A/W)\), the associated prime
ideal of \(\mathcal O_K\) has norm \(p^f\); retain it
exactly when \(p^f\le T\). It is explicitly the inverse
image of the subspace \(W\). To output an integer basis,
put a basis of \(W\) in reduced row echelon form, lift its
entries to \(\{0,\ldots,p-1\}\), and add \(p\) times each
coordinate vector at a nonpivot position. These \(n\)
vectors generate precisely the inverse-image lattice:
subtracting the lifted rows using the integral pivot
coordinates leaves zero pivot coordinates and multiples
of \(p\) elsewhere. Their determinant has absolute value
\(p^f\), as also follows from the quotient size.

All algebra dimensions are at most \(n\). The number of
scanned scalars is polynomial in \(p,n\), and \(p\le T\)
is polynomial in the input length by (C12.1). All modular
entries use \(O(\log T)\) bits. There are at most \(n\)
prime ideals above each \(p\), so the number of output
ideals is at most \(nT\), with polynomial total encoding
length. This proves the claimed deterministic bit
complexity. The generation conclusion follows from
(C09.2) because \(T\ge B_n\), or \(T\ge B'_n\) with
odd class number.
\end{proof}

The supplied maximal-order basis is essential to this
formulation. Starting only with a defining polynomial
does not supply it for free. The proof makes no complexity
claim about finding that basis, identifying relations,
or proving an odd-class-number promise.

\paragraph{NT-C13: finite Frobenius detection of ordinary class-field quotients.}
Under (C09.1), let \(E/K\) be abelian, unramified at every
finite place and split at every real place. If all primes
of norm at most \(B_n\) split completely in \(E\), then
\(E=K\). If the smaller bound \(B'_n\) is used instead,
then \(\operatorname{Gal}(E/K)\) has exponent at most two.
The correction-free bounds in (C09.3) have the same
respective conclusions.

\begin{proof}
Ordinary global class field theory gives a surjection
\[
       G\twoheadrightarrow\operatorname{Gal}(E/K)
\]
sending each prime class to its Frobenius. Complete
splitting kills the indicated small-prime subgroup.
At \(B_n\) that subgroup is \(G\), so the quotient is
trivial. At \(B'_n\) it contains \(G^2\); consequently
the quotient is a quotient of \(G/G^2\) and has exponent
at most two. The sharper-range versions follow identically.
\end{proof}

In geometric language, \(G=\operatorname{Pic}
(\operatorname{Spec}\mathcal O_K)\), and these tests concern
abelian finite \'etale arithmetic covers with the ordinary
real-place splitting condition. The finite prime-sum
certificate proves generation of a prescribed quotient
by its closed-point Frobenius elements. The enumeration
proof exposes the precise integral-algebra input needed
to make these points concrete. Run:
\texttt{GRH\_NIGHT/number\_theory}; primary inputs
\texttt{S-NT-GM2025-CG}, \texttt{S-NT-BDDF},
and \texttt{S-NT-GM2018-CG}.
